% Options for packages loaded elsewhere
\PassOptionsToPackage{unicode}{hyperref}
\PassOptionsToPackage{hyphens}{url}
\documentclass[
]{article}
\usepackage{xcolor}
\usepackage[margin=1in]{geometry}
\usepackage{amsmath,amssymb}
\setcounter{secnumdepth}{-\maxdimen} % remove section numbering
\usepackage{iftex}
\ifPDFTeX
  \usepackage[T1]{fontenc}
  \usepackage[utf8]{inputenc}
  \usepackage{textcomp} % provide euro and other symbols
\else % if luatex or xetex
  \usepackage{unicode-math} % this also loads fontspec
  \defaultfontfeatures{Scale=MatchLowercase}
  \defaultfontfeatures[\rmfamily]{Ligatures=TeX,Scale=1}
\fi
\usepackage{lmodern}
\ifPDFTeX\else
  % xetex/luatex font selection
\fi
% Use upquote if available, for straight quotes in verbatim environments
\IfFileExists{upquote.sty}{\usepackage{upquote}}{}
\IfFileExists{microtype.sty}{% use microtype if available
  \usepackage[]{microtype}
  \UseMicrotypeSet[protrusion]{basicmath} % disable protrusion for tt fonts
}{}
\makeatletter
\@ifundefined{KOMAClassName}{% if non-KOMA class
  \IfFileExists{parskip.sty}{%
    \usepackage{parskip}
  }{% else
    \setlength{\parindent}{0pt}
    \setlength{\parskip}{6pt plus 2pt minus 1pt}}
}{% if KOMA class
  \KOMAoptions{parskip=half}}
\makeatother
\usepackage{longtable,booktabs,array}
\usepackage{calc} % for calculating minipage widths
% Correct order of tables after \paragraph or \subparagraph
\usepackage{etoolbox}
\makeatletter
\patchcmd\longtable{\par}{\if@noskipsec\mbox{}\fi\par}{}{}
\makeatother
% Allow footnotes in longtable head/foot
\IfFileExists{footnotehyper.sty}{\usepackage{footnotehyper}}{\usepackage{footnote}}
\makesavenoteenv{longtable}
\usepackage{graphicx}
\makeatletter
\newsavebox\pandoc@box
\newcommand*\pandocbounded[1]{% scales image to fit in text height/width
  \sbox\pandoc@box{#1}%
  \Gscale@div\@tempa{\textheight}{\dimexpr\ht\pandoc@box+\dp\pandoc@box\relax}%
  \Gscale@div\@tempb{\linewidth}{\wd\pandoc@box}%
  \ifdim\@tempb\p@<\@tempa\p@\let\@tempa\@tempb\fi% select the smaller of both
  \ifdim\@tempa\p@<\p@\scalebox{\@tempa}{\usebox\pandoc@box}%
  \else\usebox{\pandoc@box}%
  \fi%
}
% Set default figure placement to htbp
\def\fps@figure{htbp}
\makeatother
\setlength{\emergencystretch}{3em} % prevent overfull lines
\providecommand{\tightlist}{%
  \setlength{\itemsep}{0pt}\setlength{\parskip}{0pt}}
\usepackage{bookmark}
\IfFileExists{xurl.sty}{\usepackage{xurl}}{} % add URL line breaks if available
\urlstyle{same}
\hypersetup{
  pdftitle={基于 Barra 风格因子的产品超额收益归因分析},
  pdfauthor={武宇翔},
  hidelinks,
  pdfcreator={LaTeX via pandoc}}

\title{基于 Barra 风格因子的产品超额收益归因分析}
\author{武宇翔}
\date{2026-06-19}

\begin{document}
\maketitle

\section{第一部分：Barra
因子与产品超额收益的相关性分析}\label{ux7b2cux4e00ux90e8ux5206barra-ux56e0ux5b50ux4e0eux4ea7ux54c1ux8d85ux989dux6536ux76caux7684ux76f8ux5173ux6027ux5206ux6790}

本部分为探索性分析，目的是初步筛选与产品超额收益相关的 Barra
风格因子，并识别因子之间的多重共线性，为后续 PCA 与回归归因做准备。

\subsection{1.1 读取 Barra
因子并与超额收益对齐}\label{ux8bfbux53d6-barra-ux56e0ux5b50ux5e76ux4e0eux8d85ux989dux6536ux76caux5bf9ux9f50}

Barra 数据为\textbf{日频因子组合累积指数}（首日 =
1），需转换为周收益率后再与产品超额收益计算相关系数。

\subsection{1.2
因子指数转周收益率}\label{ux56e0ux5b50ux6307ux6570ux8f6cux5468ux6536ux76caux7387}

\subsection{1.3 产品超额收益与各 Barra
因子的相关系数}\label{ux4ea7ux54c1ux8d85ux989dux6536ux76caux4e0eux5404-barra-ux56e0ux5b50ux7684ux76f8ux5173ux7cfbux6570}

\begin{longtable}[]{@{}llrrl@{}}
\caption{表 1：产品超额收益与各 Barra 因子周收益率的 Pearson 相关系数（*
p\textless0.05, ** p\textless0.01）}\tabularnewline
\toprule\noalign{}
因子中文名 & 因子 & 相关系数 & p值 & 显著性 \\
\midrule\noalign{}
\endfirsthead
\toprule\noalign{}
因子中文名 & 因子 & 相关系数 & p值 & 显著性 \\
\midrule\noalign{}
\endhead
\bottomrule\noalign{}
\endlastfoot
规模 & size & -0.5204 & 0.0000 & ** \\
杠杆 & leverage & 0.3259 & 0.0125 & * \\
中市值 & mid\_cap & -0.2666 & 0.0431 & * \\
长期反转 & long\_term\_reversal & 0.2560 & 0.0524 & . \\
成长 & growth & -0.2344 & 0.0765 & . \\
流动性 & liquidity & 0.2120 & 0.1102 & \\
盈利能力 & profitability & -0.1890 & 0.1552 & \\
盈利收益率 & earning\_yield & -0.1772 & 0.1833 & \\
波动率 & volatility & -0.1041 & 0.4370 & \\
投资质量 & investment\_quality & 0.0680 & 0.6118 & \\
股息 & dividend & -0.0627 & 0.6403 & \\
账面市值比 & book\_to\_price & 0.0280 & 0.8350 & \\
动量 & momentum & 0.0055 & 0.9676 & \\
盈利波动 & earning\_variability & -0.0048 & 0.9714 & \\
贝塔 & beta & -0.0023 & 0.9861 & \\
\end{longtable}

\subsection{1.4 Barra
因子之间的相关矩阵}\label{barra-ux56e0ux5b50ux4e4bux95f4ux7684ux76f8ux5173ux77e9ux9635}

\begin{center}\includegraphics{基于Barra风格因子的产品超额收益归因分析_files/figure-latex/part1-corr-matrix-1} \end{center}

\subsection{1.5
超额收益与主要因子的散点图}\label{ux8d85ux989dux6536ux76caux4e0eux4e3bux8981ux56e0ux5b50ux7684ux6563ux70b9ux56fe}

选取与超额收益相关性最强的 3 个因子作散点图，直观展示线性关系。

\begin{center}\includegraphics{基于Barra风格因子的产品超额收益归因分析_files/figure-latex/part1-scatter-1} \end{center}

\section{第二部分：PCA 对 Barra
因子降维}\label{ux7b2cux4e8cux90e8ux5206pca-ux5bf9-barra-ux56e0ux5b50ux964dux7ef4}

第一部分发现部分 Barra 因子之间存在中等程度相关（如 size 与
mid\_cap），若将全部 15
个因子直接放入多元回归，可能带来系数不稳定问题。本部分对 Barra
因子周收益率矩阵做 \textbf{PCA
主成分分析}，在保留主要因子波动信息的同时降低维度，并为第三部分主成分回归做准备。

\subsection{2.1
因子收益率矩阵标准化}\label{ux56e0ux5b50ux6536ux76caux7387ux77e9ux9635ux6807ux51c6ux5316}

PCA
前需对各因子收益率\textbf{标准化}（减均值、除以标准差），避免波动幅度大的因子
dominate 主成分。

\subsection{2.2 主成分分析}\label{ux4e3bux6210ux5206ux5206ux6790}

对标准化后的因子矩阵做 PCA，提取主成分得分（PC
scores）与载荷（loadings）。

\subsection{2.3 方差解释率}\label{ux65b9ux5deeux89e3ux91caux7387}

观察各主成分解释了多少因子总波动，并据此确定后续回归使用的主成分个数。

\begin{longtable}[]{@{}lrr@{}}
\caption{表 3：各主成分方差解释率（前 6 个）}\tabularnewline
\toprule\noalign{}
主成分 & 方差解释率 & 累积解释率 \\
\midrule\noalign{}
\endfirsthead
\toprule\noalign{}
主成分 & 方差解释率 & 累积解释率 \\
\midrule\noalign{}
\endhead
\bottomrule\noalign{}
\endlastfoot
PC1 & 0.2342 & 0.2342 \\
PC2 & 0.1465 & 0.3807 \\
PC3 & 0.1311 & 0.5118 \\
PC4 & 0.1041 & 0.6159 \\
PC5 & 0.0713 & 0.6871 \\
PC6 & 0.0541 & 0.7412 \\
\end{longtable}

\begin{center}\includegraphics{基于Barra风格因子的产品超额收益归因分析_files/figure-latex/part2-variance-1} \end{center}

\subsection{2.4
主成分载荷分析}\label{ux4e3bux6210ux5206ux8f7dux8377ux5206ux6790}

载荷（loading）反映各原始因子在主成分上的权重，用于解读 PC1、PC2
等的经济含义。

\begin{longtable}[]{@{}
  >{\raggedright\arraybackslash}p{(\linewidth - 12\tabcolsep) * \real{0.2532}}
  >{\raggedright\arraybackslash}p{(\linewidth - 12\tabcolsep) * \real{0.1392}}
  >{\raggedright\arraybackslash}p{(\linewidth - 12\tabcolsep) * \real{0.2532}}
  >{\raggedleft\arraybackslash}p{(\linewidth - 12\tabcolsep) * \real{0.0886}}
  >{\raggedleft\arraybackslash}p{(\linewidth - 12\tabcolsep) * \real{0.0886}}
  >{\raggedleft\arraybackslash}p{(\linewidth - 12\tabcolsep) * \real{0.0886}}
  >{\raggedleft\arraybackslash}p{(\linewidth - 12\tabcolsep) * \real{0.0886}}@{}}
\caption{表 4：前 4 个主成分的因子载荷}\tabularnewline
\toprule\noalign{}
\begin{minipage}[b]{\linewidth}\raggedright
\end{minipage} & \begin{minipage}[b]{\linewidth}\raggedright
因子中文名
\end{minipage} & \begin{minipage}[b]{\linewidth}\raggedright
因子
\end{minipage} & \begin{minipage}[b]{\linewidth}\raggedleft
PC1
\end{minipage} & \begin{minipage}[b]{\linewidth}\raggedleft
PC2
\end{minipage} & \begin{minipage}[b]{\linewidth}\raggedleft
PC3
\end{minipage} & \begin{minipage}[b]{\linewidth}\raggedleft
PC4
\end{minipage} \\
\midrule\noalign{}
\endfirsthead
\toprule\noalign{}
\begin{minipage}[b]{\linewidth}\raggedright
\end{minipage} & \begin{minipage}[b]{\linewidth}\raggedright
因子中文名
\end{minipage} & \begin{minipage}[b]{\linewidth}\raggedright
因子
\end{minipage} & \begin{minipage}[b]{\linewidth}\raggedleft
PC1
\end{minipage} & \begin{minipage}[b]{\linewidth}\raggedleft
PC2
\end{minipage} & \begin{minipage}[b]{\linewidth}\raggedleft
PC3
\end{minipage} & \begin{minipage}[b]{\linewidth}\raggedleft
PC4
\end{minipage} \\
\midrule\noalign{}
\endhead
\bottomrule\noalign{}
\endlastfoot
beta & 贝塔 & beta & 0.442 & 0.081 & 0.131 & -0.069 \\
book\_to\_price & 账面市值比 & book\_to\_price & -0.111 & -0.020 & 0.500
& 0.107 \\
dividend & 股息 & dividend & -0.222 & -0.199 & -0.090 & 0.274 \\
earning\_variability & 盈利波动 & earning\_variability & -0.087 & 0.438
& 0.052 & 0.121 \\
earning\_yield & 盈利收益率 & earning\_yield & -0.075 & 0.305 & 0.078 &
0.482 \\
growth & 成长 & growth & -0.139 & -0.221 & 0.116 & 0.482 \\
investment\_quality & 投资质量 & investment\_quality & -0.315 & -0.217 &
0.311 & 0.211 \\
leverage & 杠杆 & leverage & -0.054 & -0.234 & 0.499 & -0.065 \\
liquidity & 流动性 & liquidity & 0.166 & -0.466 & 0.082 & -0.163 \\
long\_term\_reversal & 长期反转 & long\_term\_reversal & -0.217 & -0.320
& -0.061 & -0.346 \\
mid\_cap & 中市值 & mid\_cap & 0.267 & -0.390 & -0.156 & 0.290 \\
momentum & 动量 & momentum & 0.395 & 0.060 & 0.241 & 0.022 \\
profitability & 盈利能力 & profitability & 0.395 & -0.039 & 0.067 &
0.161 \\
size & 规模 & size & 0.218 & -0.198 & -0.386 & 0.352 \\
volatility & 波动率 & volatility & -0.320 & -0.044 & -0.333 & -0.016 \\
\end{longtable}

\begin{center}\includegraphics{基于Barra风格因子的产品超额收益归因分析_files/figure-latex/part2-loadings-1} \end{center}

\subsection{2.5
主成分与产品超额收益的相关性}\label{ux4e3bux6210ux5206ux4e0eux4ea7ux54c1ux8d85ux989dux6536ux76caux7684ux76f8ux5173ux6027}

在回归之前，先看各主成分与产品超额收益的线性关系，作为第三部分的对照。

\begin{longtable}[]{@{}lrrl@{}}
\caption{表 5：产品超额收益与各主成分的 Pearson 相关系数}\tabularnewline
\toprule\noalign{}
主成分 & 相关系数 & p值 & 显著性 \\
\midrule\noalign{}
\endfirsthead
\toprule\noalign{}
主成分 & 相关系数 & p值 & 显著性 \\
\midrule\noalign{}
\endhead
\bottomrule\noalign{}
\endlastfoot
PC1 & -0.1213 & 0.3644 & \\
PC2 & -0.0305 & 0.8200 & \\
PC3 & 0.3063 & 0.0194 & * \\
PC4 & -0.5060 & 0.0001 & ** \\
\end{longtable}

\section{第三部分：多因子回归归因}\label{ux7b2cux4e09ux90e8ux5206ux591aux56e0ux5b50ux56deux5f52ux5f52ux56e0}

本部分在前两部分基础上，分别用\textbf{筛选后的原始 Barra 因子}（版本
A）和 \textbf{PCA 主成分}（版本
B）对产品超额收益做多元线性回归，比较两种方式的解释力，并判断超额收益更多来自风格暴露还是无法被因子解释的
alpha。

回归模型：

\[
\text{超额收益}_t = \alpha + \sum_i \beta_i \cdot \text{因子}_t + \varepsilon_t
\]

\subsection{3.1
产品累计超额收益走势}\label{ux4ea7ux54c1ux7d2fux8ba1ux8d85ux989dux6536ux76caux8d70ux52bf}

\begin{center}\includegraphics{基于Barra风格因子的产品超额收益归因分析_files/figure-latex/part3-cum-excess-1} \end{center}

\subsection{3.2 版本
A：筛选因子多元回归}\label{ux7248ux672c-aux7b5bux9009ux56e0ux5b50ux591aux5143ux56deux5f52}

使用第一部分筛选出的因子（\texttt{筛选因子}）作为解释变量，估计各风格暴露及截距
alpha。

\begin{longtable}[]{@{}
  >{\raggedright\arraybackslash}p{(\linewidth - 12\tabcolsep) * \real{0.1642}}
  >{\raggedright\arraybackslash}p{(\linewidth - 12\tabcolsep) * \real{0.2836}}
  >{\raggedleft\arraybackslash}p{(\linewidth - 12\tabcolsep) * \real{0.1194}}
  >{\raggedleft\arraybackslash}p{(\linewidth - 12\tabcolsep) * \real{0.1045}}
  >{\raggedleft\arraybackslash}p{(\linewidth - 12\tabcolsep) * \real{0.1194}}
  >{\raggedleft\arraybackslash}p{(\linewidth - 12\tabcolsep) * \real{0.1045}}
  >{\raggedright\arraybackslash}p{(\linewidth - 12\tabcolsep) * \real{0.1045}}@{}}
\caption{表 6：版本 A --- 筛选因子回归（size + leverage + mid\_cap +
long\_term\_reversal + growth）}\tabularnewline
\toprule\noalign{}
\begin{minipage}[b]{\linewidth}\raggedright
变量中文名
\end{minipage} & \begin{minipage}[b]{\linewidth}\raggedright
变量
\end{minipage} & \begin{minipage}[b]{\linewidth}\raggedleft
估计值
\end{minipage} & \begin{minipage}[b]{\linewidth}\raggedleft
标准误
\end{minipage} & \begin{minipage}[b]{\linewidth}\raggedleft
t值
\end{minipage} & \begin{minipage}[b]{\linewidth}\raggedleft
p值
\end{minipage} & \begin{minipage}[b]{\linewidth}\raggedright
显著性
\end{minipage} \\
\midrule\noalign{}
\endfirsthead
\toprule\noalign{}
\begin{minipage}[b]{\linewidth}\raggedright
变量中文名
\end{minipage} & \begin{minipage}[b]{\linewidth}\raggedright
变量
\end{minipage} & \begin{minipage}[b]{\linewidth}\raggedleft
估计值
\end{minipage} & \begin{minipage}[b]{\linewidth}\raggedleft
标准误
\end{minipage} & \begin{minipage}[b]{\linewidth}\raggedleft
t值
\end{minipage} & \begin{minipage}[b]{\linewidth}\raggedleft
p值
\end{minipage} & \begin{minipage}[b]{\linewidth}\raggedright
显著性
\end{minipage} \\
\midrule\noalign{}
\endhead
\bottomrule\noalign{}
\endlastfoot
截距 (α) & (Intercept) & 0.0031 & 0.0011 & 2.7523 & 0.0081 & ** \\
规模 & size & -0.7205 & 0.2413 & -2.9863 & 0.0043 & ** \\
杠杆 & leverage & 0.8868 & 0.5432 & 1.6325 & 0.1086 & \\
中市值 & mid\_cap & 0.1427 & 0.4483 & 0.3183 & 0.7515 & \\
长期反转 & long\_term\_reversal & 0.5769 & 0.3189 & 1.8091 & 0.0762 &
. \\
成长 & growth & -1.3043 & 0.6488 & -2.0103 & 0.0496 & * \\
\end{longtable}

\textbf{p 值怎么读？} p
值表示「该变量对超额收益没有解释力」时，仍观察到当前结果的概率。\textbf{p
值越小，该因子在控制其他变量后，对超额收益的解释越重要}（越不像随机噪声）。一般
p \textless{} 0.05 认为显著，p \textless{} 0.01 认为高度显著。

\textbf{本表结论（版本 A）：}

\begin{itemize}
\tightlist
\item
  \textbf{截距 α}（p =
  0.0081）：显著为正，说明扣掉所列风格因子后，产品仍有稳定的每周 alpha。
\item
  \textbf{对超额解释最重要的风格因子}（p \textless{} 0.05，按 p
  值从小到大）： 规模（β = -0.7205，p = 0.0043）； 成长（β = -1.3043，p
  = 0.0496）。
\item
  \textbf{弱显著}（0.05 ≤ p \textless{} 0.10）：长期反转（p =
  0.0762），仅供参考。
\item
  \textbf{对超额解释不显著}（p ≥ 0.10）：杠杆（p = 0.1086）、中市值（p =
  0.7515），在多元模型中可认为对样本期内超额波动\textbf{解释力有限}。
\end{itemize}

\subsection{3.3 版本
B：主成分多元回归}\label{ux7248ux672c-bux4e3bux6210ux5206ux591aux5143ux56deux5f52}

使用第二部分建议的主成分（\texttt{使用主成分}）作为解释变量，避免原始因子共线性带来的系数不稳定。

\begin{longtable}[]{@{}llrrrrl@{}}
\caption{表 7：版本 B --- 主成分回归（PC1 + PC2 + PC3 +
PC4）}\tabularnewline
\toprule\noalign{}
变量中文名 & 变量 & 估计值 & 标准误 & t值 & p值 & 显著性 \\
\midrule\noalign{}
\endfirsthead
\toprule\noalign{}
变量中文名 & 变量 & 估计值 & 标准误 & t值 & p值 & 显著性 \\
\midrule\noalign{}
\endhead
\bottomrule\noalign{}
\endlastfoot
截距 (α) & (Intercept) & 0.0044 & 0.0011 & 3.9962 & 0.0002 & ** \\
PC1 & PC1 & -0.0007 & 0.0006 & -1.1087 & 0.2726 & \\
PC2 & PC2 & -0.0002 & 0.0007 & -0.2790 & 0.7813 & \\
PC3 & PC3 & 0.0022 & 0.0008 & 2.7992 & 0.0071 & ** \\
PC4 & PC4 & -0.0041 & 0.0009 & -4.6250 & 0.0000 & ** \\
\end{longtable}

\textbf{p 值怎么读？} 与版本 A 相同：主成分系数 γ 的 p 值越小，说明该
\textbf{PC 对超额收益越重要}。PC
是多个原始因子的线性组合，显著的主成分需结合第二部分载荷表理解其风格含义。

\textbf{本表结论（版本 B）：}

\begin{itemize}
\tightlist
\item
  \textbf{截距 α}（p = 0.0002）：显著，扣除主成分后仍有无法被 PC
  解释的正向超额。
\item
  \textbf{对超额解释最重要的主成分}（p \textless{} 0.05，p
  越小越重要）： PC4（γ = -0.0041，p = 0.0000）； PC3（γ = +0.0022，p =
  0.0071）。
\item
  \textbf{对超额解释不显著}（p ≥ 0.05）：PC1（p = 0.2726）、PC2（p =
  0.7813）。这些 PC 虽解释因子间波动，但对产品超额\textbf{贡献有限}（与
  2.5 结论一致）。
\end{itemize}

\subsection{3.4
两种模型对比}\label{ux4e24ux79cdux6a21ux578bux5bf9ux6bd4}

\begin{longtable}[]{@{}lrr@{}}
\caption{表 8：两种回归模型拟合效果对比}\tabularnewline
\toprule\noalign{}
指标 & 版本A\_筛选因子 & 版本B\_主成分 \\
\midrule\noalign{}
\endfirsthead
\toprule\noalign{}
指标 & 版本A\_筛选因子 & 版本B\_主成分 \\
\midrule\noalign{}
\endhead
\bottomrule\noalign{}
\endlastfoot
R² & 0.3889 & 0.3655 \\
调整 R² & 0.3301 & 0.3176 \\
AIC & -383.7043 & -383.5310 \\
BIC & -369.2812 & -371.1683 \\
残差标准误 & 0.0083 & 0.0084 \\
\end{longtable}

\subsection{3.5 拟合值与残差：风格解释 vs
alpha}\label{ux62dfux5408ux503cux4e0eux6b8bux5deeux98ceux683cux89e3ux91ca-vs-alpha}

表 9 对比 \textbf{版本 A（筛选因子回归）} 与 \textbf{版本
B（主成分回归）} 两种模型下，样本期平均超额收益如何拆成「截距
α（风格无法解释的部分）」与「风格拟合（因子/主成分能解释的部分）」；左列、右列分别对应两种归因口径，便于对照同一指标在两种模型下的差异。

\begin{longtable}[]{@{}
  >{\raggedright\arraybackslash}p{(\linewidth - 4\tabcolsep) * \real{0.4930}}
  >{\raggedleft\arraybackslash}p{(\linewidth - 4\tabcolsep) * \real{0.2676}}
  >{\raggedleft\arraybackslash}p{(\linewidth - 4\tabcolsep) * \real{0.2394}}@{}}
\caption{表 9：超额收益归因分解（左列 = 版本 A，右列 = 版本
B）}\tabularnewline
\toprule\noalign{}
\begin{minipage}[b]{\linewidth}\raggedright
指标
\end{minipage} & \begin{minipage}[b]{\linewidth}\raggedleft
版本 A（筛选因子）
\end{minipage} & \begin{minipage}[b]{\linewidth}\raggedleft
版本 B（主成分）
\end{minipage} \\
\midrule\noalign{}
\endfirsthead
\toprule\noalign{}
\begin{minipage}[b]{\linewidth}\raggedright
指标
\end{minipage} & \begin{minipage}[b]{\linewidth}\raggedleft
版本 A（筛选因子）
\end{minipage} & \begin{minipage}[b]{\linewidth}\raggedleft
版本 B（主成分）
\end{minipage} \\
\midrule\noalign{}
\endhead
\bottomrule\noalign{}
\endlastfoot
周均超额收益 & 0.0044 & 0.0044 \\
截距 α（周均无法被因子解释的部分） & 0.0031 & 0.0044 \\
风格拟合均值 & 0.0013 & 0.0000 \\
残差均值 & -0.0000 & -0.0000 \\
\end{longtable}

\textbf{读表说明：} 每一行有\textbf{两个独立数字}，左列为版本
A（筛选因子回归），右列为版本
B（主成分回归）。例如「周均超额收益」一行：版本 A =
\textbf{0.0044}，版本 B =
\textbf{0.0044}（二者应接近，均为样本期平均超额）。

\begin{center}\includegraphics{基于Barra风格因子的产品超额收益归因分析_files/figure-latex/part3-attribution-plot-a-1} \end{center}

\begin{center}\includegraphics{基于Barra风格因子的产品超额收益归因分析_files/figure-latex/part3-attribution-b-1} \end{center}

\subsection{3.6
研究结论与局限性}\label{ux7814ux7a76ux7ed3ux8bbaux4e0eux5c40ux9650ux6027}

\subsubsection{研究结论}\label{ux7814ux7a76ux7ed3ux8bba}

在 58 周样本内，筛选 Barra 因子可解释约 \textbf{39\%} 的超额波动，调整
R² 约 \textbf{33\%}。规模、成长暴露显著为负，截距 alpha 显著为正。PCA
主成分回归调整 R² 约 \textbf{32\%}，略低于因子版，因 PCA
按因子波动而非超额相关性构造，PC1、PC2
对超额解释力弱。综合来看，产品超额与\textbf{规模风格}关系最明显，但仍有约
\textbf{61\%} 波动无法被所选 Barra
因子线性解释，需结合样本期短、归因口径简化等局限性谨慎解读。

\subsubsection{产品涨跌与什么有关？（Barra
风格解读）}\label{ux4ea7ux54c1ux6da8ux8dccux4e0eux4ec0ux4e48ux6709ux5173barra-ux98ceux683cux89e3ux8bfb}

结合表 1 相关性与表 6 回归，样本期内产品超额（相对中证
1000）主要与以下因素有关：

\begin{longtable}[]{@{}
  >{\raggedright\arraybackslash}p{(\linewidth - 4\tabcolsep) * \real{0.2727}}
  >{\raggedright\arraybackslash}p{(\linewidth - 4\tabcolsep) * \real{0.2727}}
  >{\raggedright\arraybackslash}p{(\linewidth - 4\tabcolsep) * \real{0.4545}}@{}}
\toprule\noalign{}
\begin{minipage}[b]{\linewidth}\raggedright
维度
\end{minipage} & \begin{minipage}[b]{\linewidth}\raggedright
判断
\end{minipage} & \begin{minipage}[b]{\linewidth}\raggedright
通俗理解
\end{minipage} \\
\midrule\noalign{}
\endhead
\bottomrule\noalign{}
\endlastfoot
\textbf{规模（size）} & 相关最强（r ≈ -0.52），回归显著负向 &
小盘/规模因子表现强时，产品超额往往偏弱；规模因子走弱时，超额往往更好。产品相对中证
1000 \textbf{偏大盘或低小盘暴露}，不是「小盘行情驱动超额」的类型 \\
\textbf{成长（growth）} & 多元回归显著负向 &
成长风格涨得多时，超额倾向走弱；与偏价值/低成长暴露或持仓结构有关 \\
\textbf{杠杆、中市值等} & 单因子相关有一定关系，多元回归不显著 &
对超额有弱线索，但不是主因 \\
\textbf{动量、行业类 proxy 等} & 与超额几乎无线性相关 &
样本内\textbf{难以用}「追热点、动量行情」解释超额 \\
\textbf{截距 α} & 显著为正 &
扣掉规模、成长等风格后，仍有\textbf{稳定的正超额}，更像选股、配置或
Barra 未覆盖因素带来的 alpha，而非单一风格 bet \\
\end{longtable}

\subsubsection{局限性}\label{ux5c40ux9650ux6027}

\begin{itemize}
\tightlist
\item
  样本量有限（约 58 周），统计检验功效受限；
\item
  仅考察线性关系，未控制时变暴露与滞后效应；
\item
  超额收益为算术差，与 Barra 组合收益口径可能存在差异；
\item
  相关与回归结果均为样本内拟合，不代表未来持续有效。
\end{itemize}

\end{document}
